\chapter{Results}

\section{Selected design and evaluation contract}

The primary result in this report is run
\code{07acbcab1cd346cca24eab31919496c0}, which selected
\code{pipeline\_ep0} using the PVT rank ``most robust.'' The target was
locked-phase DFE eye height $>\SI{0.1}{V}$, eye width $>0.4$ UI, decision
margin $>\SI{0}{V}$, and CTLE power $<\SI{15}{mW}$. The run used PPO v3
checkpoint \code{ppo\_v3\_tt\_1000\_policy.pt}, the IEEE/IBM 20-dB-through
development channel, port map 1,3,2,4, and candidate-fidelity PVT evaluation.

\begin{table}[htbp]
\centering
\caption{Eight parameters of the selected design used in the completed
27-condition run.}
\label{tab:selected-design}
\begin{tabular}{@{}l r l@{}}
\toprule
Parameter & Value & Unit \\
\midrule
$R_{\mathrm{LOAD}}$ & 2457.143 & $\Omega$ \\
$R_{\mathrm{DEG}}$ & 485.714 & $\Omega$ \\
$C_{\mathrm{DEG}}$ & 0.485714 & pF \\
$I_{\mathrm{TAIL}}$ & 0.670 & mA \\
Behavioral DFE tap & 57.143 & mV \\
Matched MOS width $W$ & 10.378 & $\mu$m \\
Matched MOS length $L$ & 0.150 & $\mu$m \\
MOS multiplier $M$ & 1 & integer \\
\bottomrule
\end{tabular}
\end{table}

The channel checksum is
\path{ce6002f8663623ebc0992c458c3d518faf79ccfeb65d1ee59c8c638bccd66581}.
The PVT request contains 512 bits; each corner records 425 useful simulated bits
after alignment and warm-up. These identifiers are necessary because results
from other candidates and channels are not interchangeable.

\section{Completed reduced 27-condition PVT run}

All 27 evaluated conditions passed over TT/SS/FF,
\SIlist{1.71;1.80;1.89}{V}, and \SIlist{0;27;125}{\degreeCelsius}.
\Cref{fig:pvt-heatmaps} displays the recorded values for every condition.

\begin{landscape}
\begin{figure}[p]
  \centering
  \includegraphics[width=0.98\linewidth]{pvt27-all-metrics-heatmaps.png}
  \caption{Measured performance across the completed 512-bit, 27-condition
  candidate-fidelity run. All values come from run \code{07acbcab...}; all 27
  conditions passed. Power is CTLE power. The DFE metrics are behavioral, and
  zero finite-pattern errors are not a BER result. This reduced grid excludes
  \SI{75}{\degreeCelsius}, SF, and FS.}
  \label{fig:pvt-heatmaps}
\end{figure}
\end{landscape}

\begin{figure}[htbp]
  \centering
  \includegraphics[width=0.96\textwidth]{pvt27-worst-case-spec-margins.png}
  \caption{Worst recorded corner for each evaluated specification in run
  \code{07acbcab...}. The plot reports measured values and locations directly;
  it does not normalize the margins into an aggregate score.}
  \label{fig:pvt-worst}
\end{figure}

The worst locked-phase eye height was \SI{1.16620}{V} at FF,
\SI{1.71}{V}, \SI{125}{\degreeCelsius}; worst eye width was 0.85 UI at the
same corner; and worst DFE margin was \SI{0.55852}{V}. CTLE power ranged from
\SI{1.1457}{mW} to \SI{1.2663}{mW}. Input-referred noise ranged from
\SI{0.328262}{mV RMS} to \SI{0.487290}{mV RMS}. The least favorable HD3
(largest value) was \SI{-69.6987}{dB} at FF, \SI{1.71}{V},
\SI{125}{\degreeCelsius}. Peaking remained between 3.18078 and
\SI{3.75745}{dB}.

These margins demonstrate robust behavior over the reduced grid, but not the
unmeasured competition corners or physical implementation effects.

\section{CTLE AC response}

\begin{landscape}
\begin{figure}[p]
  \centering
  \includegraphics[width=0.96\linewidth]{ctle-ac-responses-27-corners.png}
  \caption{Regenerated AC-only ngspice sweeps for the selected design and the 27
  recorded process/voltage/temperature conditions. The shaded region marks the
  intended \SIrange{1.25}{2.5}{GHz} peak band and the dashed line marks the
  \SI{2.5}{GHz} Nyquist frequency. These curves were regenerated from the run
  configuration; they are not stored transient/noise/HD3 waveforms and do not by
  themselves prove PVT qualification.}
  \label{fig:ac-corners}
\end{figure}
\end{landscape}

\section{PPO v3 training evidence}

The recorded seed-101 v3 training run completed 100 updates and 692 real-backend
evaluations. \Cref{fig:ppo-training} shows reward, cumulative satisfying steps,
and the first failed stage. The dominant AC failures explain why staged
diagnostics and validity-aware state are important. The curve is one training
history, not a confidence interval or proof of convergence.

\begin{landscape}
\begin{figure}[p]
  \centering
  \includegraphics[width=0.97\linewidth]{ppo-v3-training-overview.pdf}
  \caption{PPO v3 real-backend training evidence from
  \code{v3\_tt\_1000\_training.jsonl}. Reward v3 values must not be compared
  directly with v1/v2 reward scales. Cumulative satisfying steps are within-run
  observations, not independent trials.}
  \label{fig:ppo-training}
\end{figure}
\end{landscape}

\section{Fair optimizer comparison}

\begin{figure}[htbp]
  \centering
  \includegraphics[width=0.78\textwidth]{optimizer-comparison.pdf}
  \caption{The matched-budget, no-warm-start comparison used one seed and 20
  evaluations per method. Random Search found one feasible design; CEM and PPO
  found none. PPO was grid-local while the baselines used continuous draws, so
  the experiment is insufficient for a general optimizer ranking.}
  \label{fig:optimizer-comparison}
\end{figure}

PPO superiority is not demonstrated. The result instead reveals a practical
reachability limitation under the tested short-horizon configuration.

\section{Evaluation efficiency}

Hierarchical early rejection avoids later stages after a hard failure, while
exact-key caching avoids repeated identical evaluations. In a controlled
20-seed synthetic study, caching reduced evaluator calls from 1200 to 1085
(9.58\%) with bit-identical action and reward trajectories. Synthetic wall-clock
time fell by only 1.41\% because PPO computation dominated the cheap algebraic
evaluator. A matched real-SPICE cache timing experiment and an exhaustive
MOS/R/C/L baseline were not completed; therefore significant real design-time
speedup and near-optimality remain unproven.

\section{Recorded eye artifact}

\IfFileExists{../results/v3_tt_1000_eye.png}{%
\begin{figure}[htbp]
  \centering
  \includegraphics[width=\textwidth]{v3_tt_1000_eye.png}
  \caption{Repository-backed eye artifact. The left panel is the continuous
  transistor-level CTLE output after channel filtering; the right panel contains
  corrected sample points from the behavioral one-tap DFE. The DFE panel is not
  a transistor-level waveform, and the finite pattern is not BER evidence.}
  \label{fig:eye}
\end{figure}
}{\emph{The optional eye artifact was not found when this document was built.}}

\section{Result boundary}

The evidence establishes a working, versioned automation pipeline with real
ngspice in the loop, a completed reduced-grid PVT run, and preserved experiment
artifacts. It does not establish a complete physical receiver, general PPO
superiority, full five-corner PVT, PCIe compliance, total area, or near-optimality.
